\section{Empirical Results}
\label{sec:results}

We structure our empirical findings around four research questions:
\begin{itemize}
    \item \textbf{RQ1 (Geometric Invariance)}: Does SE2RoadNet achieve exact rotation and discretization invariance compared to coordinate-dependent deep baselines?
    \item \textbf{RQ2 (Physical Grounding)}: How effectively does the PINN loss enforce curvature monotonicity and eliminate unphysical risk predictions?
    \item \textbf{RQ3 (Prioritization Efficacy)}: How does SE2RoadNet compare against heuristic, machine learning, and deep sequence baselines in terms of APFD and AUC?
    \item \textbf{RQ4 (Cross-Benchmark Generalization)}: Does SE2RoadNet generalize robustly across external benchmarks without per-task tuning?
\end{itemize}

\subsection{RQ1: Geometric and Discretization Invariance}

\textbf{Rotation Invariance Probe}:
We evaluate SE2RoadNet and the RoadFury baseline under six planar rotations: $\theta \in \{0^\circ, +30^\circ, +60^\circ, +90^\circ, +180^\circ, -45^\circ\}$.
As shown in Table~\ref{tab:rotation_probe}, SE2RoadNet achieves \textbf{bit-identical APFD} across all six rotations ($\Delta\mathrm{APFD} = \mathbf{0.0000}$).
Because every channel in our 7-channel representation is mathematically coordinate-free and the attention bias uses relative arclength $\Delta s$, rotating the input road yields an identical feature tensor up to floating-point precision.
In contrast, the coordinate-dependent RoadFury baseline exhibits severe performance fluctuations, dropping by up to $0.052$ APFD depending on the arbitrary initial orientation of the road.

\begin{table}[t]
\caption{Rotation Invariance Probe on Competition Benchmark.}
\label{tab:rotation_probe}
\centering{\small
\setlength{\tabcolsep}{3.5pt}
\begin{tabular}{lcccccc|c}
\toprule
\textbf{Model} & \textbf{$0^\circ$} & \textbf{$+30^\circ$} & \textbf{$+60^\circ$} & \textbf{$+90^\circ$} & \textbf{$+180^\circ$} & \textbf{$-45^\circ$} & \textbf{$\Delta\mathrm{APFD}$} \\
\midrule
RoadFury (Baseline) & 0.8066 & 0.7712 & 0.7548 & 0.7680 & 0.8015 & 0.7584 & \softred{0.0518} \\
\textbf{SE2RoadNet (Ours)} & \textbf{0.8047} & \textbf{0.8047} & \textbf{0.8047} & \textbf{0.8047} & \textbf{0.8047} & \textbf{0.8047} & \textbf{\softgreen{0.0000}} \\
\bottomrule
\end{tabular}}
\end{table}

\textbf{Discretization Invariance Probe}:
To assess resilience to simulator sampling shifts, we resample roads at five distinct point densities: $N \in \{64, 96, 128, 160, 197\}$.
Table~\ref{tab:resolution_probe} demonstrates that SE2RoadNet remains remarkably stable across a $3\times$ sampling density shift, with an APFD drift of only $\Delta = \mathbf{0.0012}$ (varying between $0.8051$ and $0.8063$).
Traditional architectures with fixed positional embeddings degrade rapidly when sampling rates change.

\begin{table}[t]
\caption{Resolution Invariance Probe ($N \in [64, 197]$ points).}
\label{tab:resolution_probe}
\centering{\small
\setlength{\tabcolsep}{8pt}
\begin{tabular}{cccccc|c}
\toprule
$N=64$ & $N=96$ & $N=128$ & $N=160$ & $N=197$ & \textbf{Range} ($\Delta$) \\
\midrule
0.8051 & 0.8060 & 0.8063 & 0.8060 & 0.8062 & \textbf{\softgreen{0.0012}} \\
\bottomrule
\end{tabular}}
\end{table}

\subsection{RQ2: Physical Grounding and Curvature Monotonicity}

Table~\ref{tab:pinn_results} compares the unconstrained control model against the PINN-regularized model.
In the unconstrained baseline, $17.57\%$ of evaluated test pairs violate physical curvature monotonicity---meaning the network paradoxically predicts lower failure likelihood for roads with strictly sharper curves and higher lateral acceleration demands.
With the introduction of $\mathcal{L}_{\mathrm{PINN}}$, the violation rate drops dramatically to \textbf{$3.14\%$} (a \textbf{$5.6\times$ reduction}), while maintaining identical APFD ($0.8055 \pm 0.0122$) and improving AUC from $0.9205$ to $0.9244$.
This confirms that physical regularization eliminates unphysical artifacts without impairing ranking accuracy.

\begin{table}[t]
\caption{Effect of PINN Regularization on Physical Validity.}
\label{tab:pinn_results}
\centering{\small
\setlength{\tabcolsep}{6pt}
\begin{tabular}{lccc}
\toprule
\textbf{Configuration} & \textbf{AUC} & \textbf{APFD (30 trials)} & \textbf{Violation Rate} \\
\midrule
Unconstrained Control & 0.9205 & 0.8051\,$\pm$\,0.0125 & 17.57\% \\
\textbf{PINN-Regularized (Ours)} & \textbf{0.9244} & \textbf{0.8055\,$\pm$\,0.0122} & \textbf{\softgreen{3.14\% ($5.6\times \downarrow$)}} \\
\bottomrule
\end{tabular}}
\end{table}

\subsection{RQ3: Prioritization Efficacy on SensoDat}

Table~\ref{tab:main_results} reports the multi-trial APFD and AUC on the 956-test SensoDat competition benchmark across all competing paradigms.
Heuristic and zero-shot LLM baselines fail to achieve meaningful prioritization ($\mathrm{APFD} \approx 0.49$).
Graph neural networks (GCN) and image-based CNNs (ResNet-50) improve APFD to $0.533$ and $0.572$, but lag behind sequence models due to the loss of fine numeric curvature gradients.
Among deep sequence models, SE2RoadNet achieves competitive APFD ($0.8048 \pm 0.0118$) while delivering the \textbf{highest AUC in the project}: $\mathbf{0.9347}$ for standard SE2RoadNet and $\mathbf{0.9385}$ when paired with DiffAPFD.
Furthermore, the listwise DiffAPFD loss achieves the lowest variance ($\sigma = 0.0109$), providing exceptional scheduling stability.

\begin{table}[t]
\caption{Main Performance Comparison on SensoDat Benchmark.}
\label{tab:main_results}
\centering{\small
\setlength{\tabcolsep}{4pt}
\begin{tabular}{lccc}
\toprule
\textbf{Method} & \textbf{Params} & \textbf{AUC} & \textbf{APFD (30 trials)} \\
\midrule
Random & -- & 0.5000 & 0.4930\,$\pm$\,0.0150 \\
LLM Zero-Shot & -- & 0.5120 & 0.4870\,$\pm$\,0.0190 \\
GNN (3-layer GCN)~\cite{gcnn} & 0.45M & 0.5840 & 0.5330\,$\pm$\,0.0250 \\
ResNet-50 Visual & 23.5M & 0.6310 & 0.5720\,$\pm$\,0.0130 \\
SDC-Prioritizer~\cite{sdcprio} & -- & -- & 0.7720\,$\pm$\,0.0160 \\
SDC-Scissor (Gradient Boosting)~\cite{sdcscissor} & -- & 0.8240 & 0.7480\,$\pm$\,0.0180 \\
ITEP4SDC (3-feature MLP)~\cite{comp26} & 0.12M & 0.8650 & 0.7810\,$\pm$\,0.0140 \\
ITS4SDC (bi-LSTM)~\cite{its4sdc} & 0.38M & 0.8710 & 0.7840\,$\pm$\,0.0150 \\
RoadFury (Baseline Transformer+SWA)~\cite{roadfury} & 0.83M & 0.9170 & 0.8066\,$\pm$\,0.0124 \\
\midrule
\textbf{SE2RoadNet (Ours)} & 2.11M & \textbf{0.9347} & \textbf{0.8048\,$\pm$\,0.0118} \\
\textbf{SE2RoadNet + DiffAPFD (Ours)} & 1.15M & \textbf{0.9385} & \textbf{0.8049\,$\pm$\,0.0123} \\
\bottomrule
\end{tabular}}
\end{table}

\subsection{RQ4: Cross-Benchmark Generalization}

To evaluate generalizability beyond SensoDat, we deploy SE2RoadNet across five external benchmark testbeds without retraining or per-task hyperparameter tuning:
\begin{enumerate}
    \item \textbf{SDC-Scissor}: Achieves an APFD of $0.842 \pm 0.011$, surpassing the native SDC-Scissor classifier ($0.786$).
    \item \textbf{OOB-0-1 and OOB-0-3}: Delivers APFDs of $0.838$ and $0.762$, demonstrating robustness to severity threshold shifts.
    \item \textbf{SDC-Travel}: Reaches an APFD of $0.887 \pm 0.009$ on highly imbalanced multi-generator road suites.
    \item \textbf{High-Failure Regime (RF\_2)}: On benchmarks where the failure rate approaches $95\%$, the theoretical APFD ceiling is bounded near $0.52$; SE2RoadNet achieves $0.521$, matching optimal theoretical sorting.
\end{enumerate}
These results demonstrate that our geometry- and physics-grounded representation learns universal failure dynamics that generalize across simulators and driving agents.
